\documentclass[11pt]{article}
\usepackage[margin=1in]{geometry}
\usepackage{amsmath,amssymb,amsthm}
\usepackage[colorlinks=true,urlcolor=blue,linkcolor=blue]{hyperref}
\newtheorem{theorem}{Theorem}
\newtheorem{lemma}{Lemma}
\newcommand{\ind}{\mathbf{1}}
\title{A recurrence whose power ranks detect perfect squares}
\author{Principia Math in Progress}
\date{28 September 2026}
\begin{document}
\maketitle

\section*{Task 1: resolve MathDB 366795}
The rank of a constant-recursive sequence is the order of its minimal
eventual constant-coefficient linear recurrence. Conjecture 5.3 of
Rowland and Sistos Barr\'on \cite{RS} predicts eventual polynomiality
for general rank sequences and eventual pseudo-polynomiality for
particular rank sequences, which allow coefficient cancellations.
We disprove the latter assertion, and hence the combined conjecture.

Here pseudo-polynomiality is understood in the periodic-coefficient,
or quasipolynomial, sense: there are a positive integer $p$ and
polynomials $P_0,\ldots,P_{p-1}$ such that the sequence equals
$P_a(M)$ for all sufficiently large $M\equiv a\pmod p$.
This makes explicit the interpretation suggested by the periodic
cancellation examples immediately preceding the conjecture.
The source does not separately define the term.

\begin{theorem}\label{thm}
The integer constant-recursive sequence
\begin{equation}\label{s}
s(n)=1+2^n+3^n-6^n+n27^n+n54^n\qquad(n\ge0)
\end{equation}
has rank eight. For every integer $M\ge5$, the rank $R(M)$ of
the sequence $(s(n)^M)_{n\ge0}$ is
\begin{equation}\label{rank}
R(M)=\frac{(M+1)(3M^2+3M+2)}2
-3(M-1)\ind_{2\mid M}-6(M-2)\ind_{M\text{ is a square}}.
\end{equation}
In particular, $R$ is not eventually quasipolynomial.
\end{theorem}

\subsection*{Rank from exponential polynomials}
We use the elementary fact that, for distinct nonzero numbers
$\lambda_1,\ldots,\lambda_v$ and nonzero polynomials $A_j$,
\begin{equation}\label{rankfact}
\operatorname{rank}\left(\sum_{j=1}^v A_j(n)\lambda_j^n\right)
=\sum_{j=1}^v(\deg A_j+1).
\end{equation}
For completeness, let $E$ be the shift operator. On $A(n)\lambda^n$,
$E-\lambda$ lowers the polynomial degree by one until it annihilates
the term; $E-\mu$, with $\mu\ne\lambda$, preserves its degree and
nonzero leading coefficient. Applying products of these operators
isolates each distinct exponential root and also proves their linear
independence. Consequently an annihilating polynomial in $E$ must
contain $(E-\lambda_j)^{\deg A_j+1}$ for each $j$, and their product
suffices. The same argument applies to eventual recurrences.

Thus \eqref{s} has minimal characteristic polynomial
\[
(z-1)(z-2)(z-3)(z-6)(z-27)^2(z-54)^2,
\]
which has integer coefficients and a nonzero constant term.
All terms of $s$ are integers, so it meets the source's hypotheses.

\subsection*{Three leading coefficient types}
Introduce independent indeterminates $t,x,y$ and put
\[
F(t,x,y)=(1+x)+y(1-x)+ty^3(1+x).
\]
Then $s(n)=F(n,2^n,3^n)$. Write
\[
F(t,x,y)^M=\sum_{i,j}C_{i,j}(t)x^iy^j.
\]
Unique prime factorization makes all numbers $2^i3^j$ distinct.
Therefore \eqref{rankfact} gives
\begin{equation}\label{sumrank}
R(M)=\sum_{C_{i,j}\ne0}(\deg C_{i,j}+1).
\end{equation}
Choosing $h$ factors of the third summand, $k$ of the second, and
$M-h-k$ of the first gives the exact expansion
\begin{equation}\label{expansion}
F^M=\sum_{\substack{h,k\ge0\\h+k\le M}}
\frac{M!}{h!k!(M-h-k)!}\,
t^hy^{3h+k}(1+x)^{M-k}(1-x)^k.
\end{equation}

Write a prospective $y$-exponent uniquely as $j=3h+r$, with
$r\in\{0,1,2\}$. Its largest possible $t$-exponent is $h$.
It is feasible exactly when $h+r\le M$: using $h-d$ third
summands instead would require $r+3d$ second summands and a total
of $h+r+2d$ factors. Hence the feasible pairs are precisely
\begin{equation}\label{feasible}
r\in\{0,1,2\},\qquad 0\le h\le M-r.
\end{equation}
For each pair and each $0\le i\le M$, the top coefficient of
$C_{i,3h+r}(t)$ is a positive multinomial multiple of
\[
A_r(i):=[x^i](1+x)^{M-r}(1-x)^r.
\]

Define the auxiliary polynomials
\[
K_k(i):=[z^k](1-z)^i(1+z)^{M-i},\qquad a=M-2i.
\]
Expansion of the binomial coefficients shows
\begin{equation}\label{dual}
\binom Mr A_r(i)=\binom Mi K_r(i).
\end{equation}
Indeed, matching the summand with $b$ negative choices on each side
gives the same factorial expression. In particular,
\begin{equation}\label{first}
K_0=1,\qquad K_1=a,\qquad K_2=\frac{a^2-M}{2}.
\end{equation}
Thus the leading coefficient never vanishes for $r=0$; for $r=1$
it vanishes only at $i=M/2$, when $M$ is even; and for $r=2$
it vanishes only at
\[
i=\frac{M-\sqrt M}{2},\quad \frac{M+\sqrt M}{2},
\]
when $M$ is a square. For square $M\ge5$ these are two distinct
integers in $[0,M]$, since $M$ and $\sqrt M$ have the same parity.

\subsection*{Exact losses, including the boundary}
At a vanishing leading coefficient, a term of degree $h-1$ is
available exactly when
\begin{equation}\label{next}
h\ge1\quad\hbox{and}\quad h\le M-r-2.
\end{equation}
By \eqref{expansion} and \eqref{dual}, it is a nonzero constant
multiple of $K_{r+3}(i)$ unless that polynomial vanishes.
Differentiating the defining generating polynomial gives
\[
(k+1)K_{k+1}=aK_k-(M-k+1)K_{k-1}.
\]
Consequently,
\begin{align}
K_4(i)&=\frac{M(M-2)}8&&\text{when }a=0,\label{k4}\\
K_5(i)&=\frac{a(M-1)(M-4)}{20}&&\text{when }a^2=M.\label{k5}
\end{align}
For example, at $a^2=M$ the recurrence successively gives
$K_3=-(M-1)a/3$ and $K_4=-M(M-1)/12$, followed by \eqref{k5}.
The quantities in \eqref{k4}--\eqref{k5} are nonzero in the
applicable cases for $M\ge5$. Thus under \eqref{next}, a vanished
leading term reduces the rank contribution from $h+1$ to $h$,
an exact loss of one.

When $h=0$, a vanished constant coefficient also loses one rank
unit. At the final two feasible values $h=M-r-1,M-r$, no smaller
$t$-exponent is possible at all: replacing $h$ by $h-d$, $d\ge1$,
would require $h+r+2d>M$ factors. The whole polynomial is then
zero, so its losses are respectively $M-r$ and $M-r+1$.
For each vanishing index $i$, the aggregate loss over all feasible
$h$ is therefore
\begin{equation}\label{loss}
\underbrace{(M-r-1)}_{h=0,\ldots,M-r-2}
+(M-r)+(M-r+1)=3(M-r).
\end{equation}
There is one such index for $r=1$ at even $M$, and two for $r=2$
at square $M$. These are disjoint $y$-exponent classes, so their
losses add even when $M$ is an even square.

If no leading terms vanished, \eqref{sumrank} would equal
\begin{align*}
(M+1)\sum_{r=0}^2\sum_{h=0}^{M-r}(h+1)
&=\frac{M+1}{2}\bigl((M+1)(M+2)+M(M+1)+(M-1)M\bigr)\\
&=\frac{(M+1)(3M^2+3M+2)}2.
\end{align*}
Subtracting exactly the losses in \eqref{loss} proves \eqref{rank}.

\subsection*{Why no eventual quasipolynomial can agree}
If $R$ were eventually quasipolynomial, subtracting its polynomial
and parity terms would make
\[
D(M)=6(M-2)\ind_{M\text{ is a square}}
\]
eventually quasipolynomial, with some positive period $p$.
On the residue class $0\pmod p$, its constituent polynomial would
vanish at all sufficiently large integers $(pt)^2+p$: these lie
strictly between $(pt)^2$ and $(pt+1)^2$, and hence are nonsquares.
That polynomial would have infinitely many roots, so would be zero.
But on the same residue class it must take the nonzero value
$6((pt)^2-2)$ at all sufficiently large squares $(pt)^2$.
This contradiction completes Theorem~\ref{thm}.

\paragraph{Scope and verification.}
All 23 pages of the current manuscript were read, and the separate
23-page author copy was fully reconciled. No journal version was
identified. An existing MathDB submission under problem 366796
claims the general, no-cancellation case and expressly leaves the
particular case open; our argument addresses that remaining case.
The accompanying verifier expands the six-term polynomial directly,
checks ranks with exact integer coefficients, and verifies the
auxiliary identities symbolically. These finite checks support the
all-$M$ proof above. Source and claim-search limits appear in the
separate audit. This is author-agent review, not independent peer review.

\section*{Tasks 2--5: proposed follow-up problems}
These are separate research proposals, not completed results.
\begin{enumerate}
\setcounter{enumi}{1}
\item Determine the minimum rank of an integer constant-recursive
sequence whose power-rank sequence is not eventually quasipolynomial.
\item Decide whether every integer constant-recursive sequence with
only simple characteristic roots has an eventually quasipolynomial
power-rank sequence.
\item Determine the power-rank sequence of
$1+2^n+3^n-6^n$ by classifying the integer zeros in the corresponding
Krawtchouk coefficient arrays.
\item For each fixed $d\ge3$, construct an integer constant-recursive
sequence with an explicit rank formula containing a nonzero correction
supported exactly on the perfect $d$th powers, or prove an obstruction
to such a construction.
\end{enumerate}

\begin{thebibliography}{9}
\bibitem{RS} Eric Rowland and Jes\'us Sistos Barr\'on,
\emph{Complexity of powers of a constant-recursive sequence},
\href{https://arxiv.org/abs/2501.14643}{arXiv:2501.14643v1} (2025),
Conjecture 5.3, p.~15.
\end{thebibliography}
\end{document}
